\section{Nonetheless, on the Day of Judgment}\label{sec:general}

If the soul's case is closed at death and its sentence carried out at once, the Christian may well ask what remains for the last day. The profession of faith received at Lyons and the constitution of Benedict~XII, having defined the immediacy of the first judgment, both go on at once to confess the second.

The profession of faith of the emperor Michael Palaeologus, received at the Second Council of Lyons in 1274, first teaches that the souls of the just, purified if they need it, are received into heaven \latin{mox}, straightway, and that the souls of those who die in mortal sin go down \latin{mox} into hell. Then it continues: \latin{Eadem sacrosancta Ecclesia Romana firmiter credit et firmiter asseverat, quod nihilominus in die iudicii omnes homines ante tribunal Christi cum suis corporibus comparebunt, reddituri de propriis factis rationem}. The same holy Roman Church firmly believes and firmly asserts that \emph{nonetheless} on the day of judgment all men will appear with their bodies before the tribunal of Christ, to render an account of their own deeds.\endnote{Second Council of Lyons (1274), profession of faith of Michael Palaeologus, DS 857--859.} Some sixty years later \latin{Benedictus Deus}, having defined the immediate vision of the blessed and the immediate descent of the damned, uses the same word again: \latin{et quod nihilominus in die iudicii omnes homines ``ante tribunal Christi'' cum suis corporibus comparebunt, reddituri de factis propriis rationem, ``ut referat unusquisque propria corporis, prout gessit, sive bonum sive malum''}.\endnote{\emph{Benedictus Deus}, DS 1002.} The words it takes from Paul are those of the second letter to the Corinthians: ``For we must all be manifested before the judgment seat of Christ, that every one may receive the proper things of the body, according as he hath done, whether it be good or evil'' (2~Cor 5:10).

\latin{Nihilominus} means nonetheless, none the less for that: the souls are already where they will be for ever, and all will still stand before Christ.

\subsection{The body that must return}

The general judgment has a condition that the particular judgment lacks: the resurrection. The Fourth Lateran Council confesses that Christ will come at the end of the age to judge the living and the dead and to render to each according to his works, to the reprobate and to the elect alike, \latin{qui omnes cum suis propriis resurgent corporibus, quae nunc gestant}---who shall all rise with their own bodies, the bodies they now bear.\endnote{Fourth Lateran Council (1215), constitution 1, \emph{Firmiter credimus}, DS 801.} The Catechism states the order: ``The resurrection of all the dead, `of both the just and the unjust,' will precede the Last Judgment.''\endnote{\emph{Catechism of the Catholic Church} 1038, official English.} On All Souls' Day the first Mass of the dead in the 1962 Missal reads the Gospel in which Christ announces it: ``the hour cometh wherein all that are in the graves shall hear the voice of the Son of God. And they that have done good things shall come forth unto the resurrection of life: but they that have done evil, unto the resurrection of judgment'' (John 5:28--29).

At death, then, what is judged is the soul, separated from the body it animated. At the last day it is the whole man who stands before Christ, soul and body together, the same man who lived and acted and sinned and repented in that flesh. Aquinas gives this as one reason for the second judgment. The soul, he says, once separated from the body, ``enters into an unchangeable condition, and receives its judgment. But the body remains subject to change down to the close of time: and therefore it must receive its reward or punishment then, in the last Judgment.''\endnote{\emph{Summa theologiae} III, q.~59, a.~5, ad~3 (English Dominican translation).} Whether it was buried with honour or left without burial, laid in consecrated ground or scattered, the body belongs to the man, and in the resurrection it too will receive its due.

The 1979 letter of the Congregation for the Doctrine of the Faith states the order of these mysteries in the language of the Council: the Church ``looks for `the glorious manifestation of our Lord, Jesus Christ' \ldots\ believing it to be distinct and deferred with respect to the situation of people immediately after death.''\endnote{Sacred Congregation for the Doctrine of the Faith, \emph{Letter on Certain Questions concerning Eschatology} (17 May 1979), point~5, official English; the internal quotation is from \emph{Dei Verbum} 4.} The general judgment takes place at that coming, which the Church still awaits.

\subsection{What a life leaves behind}

Aquinas asks why there must be a judgment of all things together, when each soul has already received its due, and he answers from the nature of judgment itself. ``Judgment cannot be passed perfectly upon any changeable subject before its consummation: just as judgment cannot be given perfectly regarding the quality of any action before its completion in itself and in its results: because many actions appear to be profitable, which in their effects prove to be hurtful.'' A man cannot be judged before his life is over, for he may still change; that is why ``It is appointed unto men once to die, and after this the Judgment.'' But a life does not wholly end when the man dies. ``Although man's temporal life in itself ends with death,'' Aquinas continues, ``still it continues dependent in a measure on what comes after it in the future.'' And he names five ways:

\begin{quote}
In one way, as it still lives on in men's memories, in which sometimes, contrary to the truth, good or evil reputations linger on. In another way in a man's children, who are so to speak something of their parent \ldots\ And yet many good men have wicked sons, and conversely. Thirdly, as to the result of his actions: just as from the deceit of Arius and other false leaders unbelief continues to flourish down to the close of the world; and even until then faith will continue to derive its progress from the preaching of the apostles. In a fourth way, as to the body, which is sometimes buried with honor and sometimes left unburied, and finally falls to dust utterly. In a fifth way, as to the things upon which a man's heart is set, such as temporal concerns, for example, some of which quickly lapse, while others endure longer.\endnote{\emph{Summa theologiae} III, q.~59, a.~5, corpus (English Dominican translation).}
\end{quote}

He concludes: ``Now all these things are submitted to the verdict of the Divine Judgment; and consequently, a perfect and public Judgment cannot be made of all these things during the course of this present time. Wherefore, there must be a final Judgment at the last day, in which everything concerning every man in every respect shall be perfectly and publicly judged.''

The third way, the result of a man's actions, is the one Aquinas illustrates from the history of the Church. Arius did not cease to deceive when he died; the deception he set going went on doing its work in souls he never met. The apostles did not cease to preach when they were martyred; the faith they planted goes on bearing fruit in every generation, and will until the end. Aquinas is careful to add that by such things ``a man neither merits nor demerits''---his merit was fixed at death---``still in a measure they accompany his reward or punishment.''\endnote{\emph{Summa theologiae} III, q.~59, a.~5, ad~1.}

Aquinas also takes up the difficulty that a second judgment would try the same case twice. To the objection drawn from Nahum in the Septuagint, ``God shall not judge twice the same thing,'' he answers with a gloss on the words themselves, ``i.e.\ in the same respect; but it is not unseemly for God to judge twice according to different respects.''\endnote{\emph{Summa theologiae} III, q.~59, a.~5, ad~2.} The Supplement to the \emph{Summa} states the principle in a single sentence: ``Each man is both an individual person and a part of the whole human race: wherefore a twofold judgment is due to him.''\endnote{\emph{Summa theologiae}, Supplement, q.~88, a.~1, ad~1, trans.\ Fathers of the English Dominican Province (London: Burns Oates and Washbourne, 1922).} The same article adds that the reward of the just is completed at the last day, for they ``will receive an increase of reward after the judgment, both from the added glory of the body and from the completion of the number of the saints.''

\subsection{The hidden justice made plain}

Augustine, beginning his book on the last judgment in \emph{The City of God}, reminds his reader that God is not idle until that day. He has judged from the beginning, and he judges now. Yet the day is rightly called the day of judgment, ``because in it there shall be no room left for the ignorant questioning why this wicked person is happy and that righteous man unhappy.''\endnote{Augustine, \emph{De civitate Dei} XX.1, trans.\ Marcus Dods, \emph{The City of God}, vol.~II (Edinburgh: T. \&\ T. Clark, 1871).} In this life, he writes, ``we do not know by what judgment of God this good man is poor and that bad man rich,'' why the innocent is condemned at the bar and the guilty goes free, why the useful man dies young and the wicked lives long. ``But when we shall have come to that judgment \ldots\ we shall then recognise the justice of all God's judgments, not only of such as shall then be pronounced, but of all which take effect from the beginning,'' even of those that now ``defy the scrutiny of human sense or insight.'' Until then, he adds, ``it is not concealed from pious minds that what is concealed is just.''\endnote{Augustine, \emph{De civitate Dei} XX.2 (Dods).}

The Supplement gives the same reason: at the general judgment ``the Divine justice will be made manifest in all things, whereas now it remains hidden.''\endnote{Supplement, q.~88, a.~1, corpus (English Dominican translation, 1922).} The Catechism describes what will then be known: ``Then through his Son Jesus Christ he will pronounce the final word on all history. We shall know the ultimate meaning of the whole work of creation and of the entire economy of salvation and understand the marvellous ways by which his Providence led everything towards its final end.''\endnote{\emph{Catechism of the Catholic Church} 1040, official English.}

Benedict~XVI, reflecting on the long cry of the poor and the innocent through history, writes that ``faith in the Last Judgement is first and foremost hope.'' ``There is a resurrection of the flesh. There is justice. There is an `undoing' of past suffering, a reparation that sets things aright.''\endnote{\emph{Spe salvi} 43, official English.} Every martyr unavenged, every slandered saint, every child wronged in secret, every good work done unseen and forgotten, waits for that day. Paul forbids the Christian to anticipate it: ``Therefore, judge not before the time: until the Lord come, who both will bring to light the hidden things of darkness and will make manifest the counsels of the hearts.'' And he ends the sentence with a promise: ``And then shall every man have praise from God'' (1~Cor 4:5).
